% GENERATED FILE. Edit canonical lessons, then run npm run generate:books.
% canonical-source-hash: fnv1a64:7e050721a9ba8fce
% canonical-lessons: SA-C466-medha, SA-C466-muhyati, SA-C466-unmadah, SA-C466-balisah, SA-C466-bodhah

\chapter{Sense and Folly}
\label{ch:sa-sense-and-folly}
\hlchaptermodality{466}

\begin{chapteropening}
\textbf{By the end of this chapter:} \emph{I can say whether someone acted wisely.}

Complete the last lesson of chapter 466: I can say whether someone acted wisely.
\end{chapteropening}

\section[\texorpdfstring{\sk{मेधा} (medhā)}{मेधा (medhā)}]{\sk{मेधा} (medhā) — mental vigour}

\label{lesson:SA-C466-medha}



\begin{quote}
Before the new one: say the Sanskrit for attention, then the Sanskrit for industrious.
\end{quote}

\subsection*{You'll want to know: \sk{मेधा}}
\textbf{\sk{मेधा}} — \emph{medhā} — \textquotedblleft{}mental vigour\textquotedblright{}.

\begin{quote}
\sk{बालस्य} \sk{मेधा} \sk{तीक्ष्णा} \sk{अस्ति।}
\end{quote}

\begin{quote}
\emph{bālasya medhā tīkṣṇā asti.}
\end{quote}

\begin{quote}
The boy has a sharp mind.
\end{quote}

A feminine noun.

\begin{cousinweb}
\textbf{\sk{मेधा}} (\emph{medhā}) MW gives \textquotedblleft{}mental vigour or power, intelligence, prudence, wisdom\textquotedblright{}; MW gives no derivation, so its origin is uncertain (MW).
\end{cousinweb}

\begin{grammarlens}[title={What you've built}]
One more word for describing what you lived and felt.
\end{grammarlens}

\subsection*{Your turn}
\begin{itemize}
\raggedright
  \item \emph{Say it:} \emph{medhā}
  \item \emph{Say it:} \emph{medhā}, once more
  \item \emph{Say it:} \emph{bālasya medhā tīkṣṇā asti.} — The boy has a sharp mind.
  \item \emph{Recall:} say the Sanskrit for attention, then the Sanskrit for industrious, then say \emph{medhā} again
\end{itemize}

\begin{tcolorbox}[breakable,colback=teal!4,colframe=teal!35!black,title={Before you move on}]
What does \sk{मेधा} mean? (\textquotedblleft{}Mental vigour\textquotedblright{}.) Say it once more.
\end{tcolorbox}
\section[\texorpdfstring{\sk{मुह्यति} (muhyati)}{मुह्यति (muhyati)}]{\sk{मुह्यति} (muhyati) — is bewildered}

\label{lesson:SA-C466-muhyati}



\begin{quote}
Before the new one: say the Sanskrit for industrious, then the Sanskrit for mental vigour.
\end{quote}

\subsection*{You'll want to know: \sk{मुह्यति}}
\textbf{\sk{मुह्यति}} — \emph{muhyati} — \textquotedblleft{}is bewildered\textquotedblright{}.

\begin{quote}
\sk{कठिनेन} \sk{प्रश्नेन} \sk{छात्रः} \sk{मुह्यति।}
\end{quote}

\begin{quote}
\emph{kaṭhinena praśnena chātraḥ muhyati.}
\end{quote}

\begin{quote}
The student is bewildered by the hard question.
\end{quote}

Past (he …): \sk{सः} \sk{मुग्धवान्}, \emph{saḥ mugdhavān}.

\begin{cousinweb}
\textbf{\sk{मुह्यति}} (\emph{muhyati}) comes from the root √muh, \textquotedblleft{}to become stupefied, to be bewildered\textquotedblright{}, the root of \textbf{\sk{मोहः}} (MW).
\end{cousinweb}

\begin{grammarlens}[title={What you've built}]
One more word for describing what you lived and felt.
\end{grammarlens}

\subsection*{Your turn}
\begin{itemize}
\raggedright
  \item \emph{Say it:} \emph{muhyati}
  \item \emph{Say it:} \emph{muhyati}, once more
  \item \emph{Say it:} \emph{kaṭhinena praśnena chātraḥ muhyati.} — The student is bewildered by the hard question.
  \item \emph{Recall:} say the Sanskrit for industrious, then the Sanskrit for mental vigour, then say \emph{muhyati} again
\end{itemize}

\begin{tcolorbox}[breakable,colback=teal!4,colframe=teal!35!black,title={Before you move on}]
What does \sk{मुह्यति} mean? (\textquotedblleft{}Is bewildered\textquotedblright{}.) Say it once more.
\end{tcolorbox}
\section[\texorpdfstring{\sk{उन्मादः} (unmādaḥ)}{उन्मादः (unmādaḥ)}]{\sk{उन्मादः} (unmādaḥ) — madness; frenzy}

\label{lesson:SA-C466-unmadah}



\begin{quote}
Before the new one: say the Sanskrit for mental vigour, then the Sanskrit for is bewildered.
\end{quote}

\subsection*{You'll want to know: \sk{उन्मादः}}
\textbf{\sk{उन्मादः}} — \emph{unmādaḥ} — \textquotedblleft{}madness; frenzy\textquotedblright{}.

\begin{quote}
\sk{उन्मादेन} \sk{सः} \sk{मार्गे} \sk{धावति।}
\end{quote}

\begin{quote}
\emph{unmādena saḥ mārge dhāvati.}
\end{quote}

\begin{quote}
In a frenzy he runs along the road.
\end{quote}

A masculine noun.

\begin{cousinweb}
\textbf{\sk{उन्मादः}} (\emph{unmādaḥ}) is built from \emph{ud}, \textquotedblleft{}up\textquotedblright{}, and the root √mad, \textquotedblleft{}to be drunk, to rejoice\textquotedblright{}; MW gives \textquotedblleft{}insanity, madness\textquotedblright{}, \textquotedblleft{}mania (as illness)\textquotedblright{} and \textquotedblleft{}intoxication\textquotedblright{} (MW).
\end{cousinweb}

\begin{grammarlens}[title={What you've built}]
One more word for describing what you lived and felt.
\end{grammarlens}

\subsection*{Your turn}
\begin{itemize}
\raggedright
  \item \emph{Say it:} \emph{unmādaḥ}
  \item \emph{Say it:} \emph{unmādaḥ}, once more
  \item \emph{Say it:} \emph{unmādena saḥ mārge dhāvati.} — In a frenzy he runs along the road.
  \item \emph{Recall:} say the Sanskrit for mental vigour, then the Sanskrit for is bewildered, then say \emph{unmādaḥ} again
\end{itemize}

\begin{tcolorbox}[breakable,colback=teal!4,colframe=teal!35!black,title={Before you move on}]
What does \sk{उन्मादः} mean? (\textquotedblleft{}Madness; frenzy\textquotedblright{}.) Say it once more.
\end{tcolorbox}
\section[\texorpdfstring{\sk{बालिशः} (bāliśaḥ)}{बालिशः (bāliśaḥ)}]{\sk{बालिशः} (bāliśaḥ) — childish; foolish}

\label{lesson:SA-C466-balisah}



\begin{quote}
Before the new one: say the Sanskrit for is bewildered, then the Sanskrit for madness.
\end{quote}

\subsection*{You'll want to know: \sk{बालिशः}}
\textbf{\sk{बालिशः}} — \emph{bāliśaḥ} — \textquotedblleft{}childish; foolish\textquotedblright{}.

\begin{quote}
\sk{वृद्धः} \sk{अपि} \sk{सः} \sk{बालिशः} \sk{अस्ति।}
\end{quote}

\begin{quote}
\emph{vṛddhaḥ api saḥ bāliśaḥ asti.}
\end{quote}

\begin{quote}
Though old, he is childish.
\end{quote}

\begin{cousinweb}
\textbf{\sk{बालिशः}} (\emph{bāliśaḥ}) MW gives \textquotedblleft{}young, childish, puerile, ignorant, simple, foolish\textquotedblright{} and \textquotedblleft{}a fool, a simpleton\textquotedblright{}; MW gives no derivation, so its origin is uncertain (MW).
\end{cousinweb}

\begin{grammarlens}[title={What you've built}]
One more word for describing what you lived and felt.
\end{grammarlens}

\subsection*{Your turn}
\begin{itemize}
\raggedright
  \item \emph{Say it:} \emph{bāliśaḥ}
  \item \emph{Say it:} \emph{bāliśaḥ}, once more
  \item \emph{Say it:} \emph{vṛddhaḥ api saḥ bāliśaḥ asti.} — Though old, he is childish.
  \item \emph{Recall:} say the Sanskrit for is bewildered, then the Sanskrit for madness, then say \emph{bāliśaḥ} again
\end{itemize}

\begin{tcolorbox}[breakable,colback=teal!4,colframe=teal!35!black,title={Before you move on}]
What does \sk{बालिशः} mean? (\textquotedblleft{}Childish; foolish\textquotedblright{}.) Say it once more.
\end{tcolorbox}
\section[\texorpdfstring{\sk{बोधः} (bodhaḥ)}{बोधः (bodhaḥ)}]{\sk{बोधः} (bodhaḥ) — understanding}

\label{lesson:SA-C466-bodhah}



\begin{quote}
Before the new one: say the Sanskrit for madness, then the Sanskrit for childish.
\end{quote}

\subsection*{You'll want to know: \sk{बोधः}}
\textbf{\sk{बोधः}} — \emph{bodhaḥ} — \textquotedblleft{}understanding\textquotedblright{}.

\begin{quote}
\sk{अनुभवेन} \sk{बोधः} \sk{वर्धते।}
\end{quote}

\begin{quote}
\emph{anubhavena bodhaḥ vardhate.}
\end{quote}

\begin{quote}
Understanding grows with experience.
\end{quote}

A masculine noun.

\begin{cousinweb}
\textbf{\sk{बोधः}} (\emph{bodhaḥ}) comes from the root √budh, \textquotedblleft{}to wake, to know\textquotedblright{}, the root of \textbf{\sk{बुद्धिः}} and \textbf{\sk{बुध्यते}}; MW gives \textquotedblleft{}knowing, understanding\textquotedblright{} and \textquotedblleft{}waking\textquotedblright{} (MW).
\end{cousinweb}

\begin{grammarlens}[title={What you've built}]
One more word for describing what you lived and felt.
\end{grammarlens}

\subsection*{Your turn}
\begin{itemize}
\raggedright
  \item \emph{Say it:} \emph{bodhaḥ}
  \item \emph{Say it:} \emph{bodhaḥ}, once more
  \item \emph{Say it:} \emph{anubhavena bodhaḥ vardhate.} — Understanding grows with experience.
  \item \emph{Recall:} say the Sanskrit for madness, then the Sanskrit for childish, then say \emph{bodhaḥ} again
\end{itemize}

\begin{tcolorbox}[breakable,colback=teal!4,colframe=teal!35!black,title={Before you move on}]
What does \sk{बोधः} mean? (\textquotedblleft{}Understanding\textquotedblright{}.) Say it once more.
\end{tcolorbox}
